\subsection{CONJUNCTION}

Let A, B be assemblies. The assembly

\[
\text { not } ((\text { not } A) \text { or } (\text { not } B))
\]

will be denoted by ``A and B''.

CS6. Let A, B, T be assemblies and x a letter. Then the assembly

\[
(T | \boldsymbol {x}) (\boldsymbol {A} \text { and } \boldsymbol {B})
\]

is identical with ``(T\textbar x)A and (T\textbar x)B''.

This is an immediate consequence of CS5 (§ 1, no. 2).

CF9. If A, B are relations in C, then ``A and B'' is a relation in C (called the conjunction of A and B).

This follows immediately from CF1 and CF2 (§ 1, no. 4).

C20. If A, B are theorems in T, then ``A and B'' is a theorem in T.

Suppose that ``A and B'' is false, that is to say,

\[
\text { not   not   } ((\text { not   } A) \text {   or   } (\text { not   } B))
\]

is true. By C16, ``(not A) or (not B)'', that is to say, \(A \Longrightarrow (\text{not } B)\) is true, hence ``not B'' is true; but this is absurd. Hence ``A and B'' is true.

C21. If A, B are relations in C, then

\[
(A \text {   and   } B) \Longrightarrow A, \quad (A \text {   and   } B) \Longrightarrow B
\]

are theorems in \(\mathcal{C}\) .

The relations (not A) \(\Longrightarrow\) ((not A) or (not B)), (not B) \(\Longrightarrow\) ((not A) or (not B)) are theorems in C, by S2 (no. 1) and C7 (no. 2). Now ((not A) or (not B)) \(\Longrightarrow\) (not (A and B)) is a theorem in C by C11. Hence (not A) \(\Longrightarrow\) (not (A and B)), (not B) \(\Longrightarrow\) (not (A and B)) are theorems in C. The result follows by applying C17.

We shall denote by ``A and B and C'' (resp. ``A or B or C'') the relation ``A and (B and C)'' (resp. ``A or (B or C)''). More generally, if

\[
A _ {1}, A _ {2}, \dots , A _ {n}
\]

are relations, we denote by `` \(A_{1}\) and \(A_{2}\) , and \ldots{} and \(A_{p}\) '' a relation which is constructed step by step by means of the convention that `` \(A_{1}\) and \(A_{2}\) and \ldots{} and \(A_{h}\) '' denotes the same relation as `` \(A_{1}\) and ( \(A_{2}\) and \ldots{} and \(A_{h}\) )''. The relation `` \(A_{1}\) or \(A_{2}\) or \ldots{} or \(A_{h}\) '' is defined similarly. The relation `` \(A_{1}\) and \(A_{2}\) and \ldots{} and \(A_{h}\) '' is a theorem in \(\mathfrak{C}\) if and only if each of the relations \(A_{1}, A_{2}, \ldots, A_{h}\) is a theorem in \(\mathfrak{C}\) .

It follows that every logical theory \(\mathcal{C}\) is equivalent to a logical theory \(\mathcal{C}'\) which has at most one explicit axiom. This is clear if \(\mathcal{C}\) has no explicit axiom. If \(\mathcal{C}\) has explicit axioms \(A_{1}, A_{2}, \ldots, A_{h}\) , let \(\mathcal{C}'\) be the theory which has the same signs and schemes as \(\mathcal{C}\) , and the explicit axiom `` \(A_{1}\) and \(A_{2}\) and \ldots{} and \(A_{h}\) ''. It is immediately seen that every axiom of \(\mathcal{C}\) (resp. \(\mathcal{C}'\) ) is a theorem of \(\mathcal{C}'\) (resp. \(\mathcal{C}\) ).

Let \(\mathfrak{C}_0\) be the theory with no explicit axioms which has the same signs as \(\mathfrak{C}\) and S1, S2, S3, S4 as its only schemes. Then the study of \(\mathfrak{C}\) reduces, in principle, to the study of \(\mathfrak{C}_0\) : for the relation \(A\) to be a theorem in \(\mathfrak{C}\) it is necessary and sufficient that there exist axioms \(A_1, A_2, \ldots, A_h\) of \(\mathfrak{C}\) such that \((A_1\) and \(A_2\) and \(\ldots\) and \(A_h) \Longrightarrow A\) is a theorem in \(\mathfrak{C}_0\) . This condition is evidently sufficient. Suppose conversely that \(A\) is a theorem in \(\mathfrak{C}\) , and let \(A_1, A_2, \ldots, A_h\) be the axioms of \(\mathfrak{C}\) appearing in a proof in \(\mathfrak{C}\) which contains \(A\) . Let \(\mathfrak{C}'\) (resp. \(\mathfrak{C}''\) ) be the theory constructed from \(\mathfrak{C}_0\) by adjoining the axioms \(A_1, A_2, \ldots, A_h\) (resp. the axiom `` \(A_1\) and \(A_2\) and \(\ldots\) and \(A_h\) ''). The proof of \(A\) in \(\mathfrak{C}\) is a proof of \(A\) in \(\mathfrak{C}'\) , therefore \(A\) is a theorem in \(\mathfrak{C}'\) and consequently in \(\mathfrak{C}''\) , because (as we noted above) \(\mathfrak{C}'\) and \(\mathfrak{C}''\) are equivalent. By the criterion of deduction, \((A_1\) and \(A_2\) and \(\ldots\) and \(A_h) \Longrightarrow A\) is a theorem in \(\mathfrak{C}_0\) . If \(\mathfrak{C}\) is contradictory, then it follows from what has been said that there exists a conjunction \(A\) of axioms of \(\mathfrak{C}\) and a relation \(R\) in \(\mathfrak{C}\) such that \(A \Longrightarrow (R \text{ and } (\text{not } R))\) is a theorem in \(\mathfrak{C}_0\) . Therefore

\[
((\text { not } R) \text { or } (\text { not   not } R)) \Longrightarrow (\text { not } A)
\]

is a theorem in \(\mathfrak{G}_0\) , and since ``(not \(R\) ) or (not not \(R\) )'' is a theorem in \(\mathfrak{G}_0\) , ``not \(A\) '' is a theorem in \(\mathfrak{G}_0\) . Conversely, if there exists a conjunction \(A\) of axioms of \(\mathfrak{G}\) such that ``not \(A\) '' is a theorem in \(\mathfrak{G}_0\) , then \(A\) and ``not \(A\) '' are theorems in \(\mathfrak{G}\) , so that \(\mathfrak{G}\) is contradictory.

